\documentclass[10pt]{article}
\usepackage[margin=0.45in,top=0.38in,bottom=0.40in]{geometry}
\usepackage{amsmath,amssymb,booktabs,graphicx,microtype,hyperref}
\usepackage[T1]{fontenc}
\usepackage{mathptmx}
\hypersetup{colorlinks=true,urlcolor=blue,citecolor=blue,linkcolor=blue}
\setlength{\parindent}{0.65em}
\setlength{\parskip}{1.4pt}
\setlength{\columnsep}{0.20in}
\setlength{\emergencystretch}{2em}
\pagestyle{plain}
\newcommand{\figcap}[2]{\par\vspace{1pt}\noindent{\footnotesize\textbf{#1} #2}\par\vspace{3pt}}
\begin{document}
\begin{center}
{\Large\bfseries SentinelRL: Budgeted Adaptive Cyber Defence in a Partially Observed Network}\par
\vspace{2pt}
{\small \textbf{Team:} Rohit Mourya (B23ES1029), Jyotiraditya Singh Rathore (B23ES1017) \quad | \quad \textbf{GitHub:} \href{https://github.com/rohitkumarmourya-maker/SentinelRL}{\texttt{github.com/rohitkumarmourya-maker/SentinelRL}}}\par
\vspace{3pt}
\end{center}
\noindent{\small\textbf{Abstract.} Enterprise network defenders must contain lateral attack propagation under noisy intrusion alerts and strict operational response budgets. We formulate adaptive cyber defense as a six-host Partially Observable Markov Decision Process (POMDP) over a 48-step horizon with a hard budget of six interventions across inspection, patching, isolation, and decoy deployment. We design a Topological-Bayesian Graph-Dueling Double-DQN combining normalized graph convolution, 3-step returns, stealth-transition belief priors, and a verification safety shield. Across 30 held-out scenarios and out-of-distribution attack-rate shifts, our agent outperforms all heuristic and tabular RL baselines in episode return ($-0.360$ vs.\ $-2.573$), critical-host compromises ($0.111$ vs.\ $0.233$), uncontained breaches ($0.044$ vs.\ $0.433$), detection delay ($0.922$ steps), and achieves zero false isolations ($0.000$).}\par
\vspace{2pt}\hrule\vspace{3pt}

\noindent\begin{minipage}[t]{0.484\textwidth}
\small
\section*{1. Problem and Environment Dynamics}
Enterprise Security Operations Centers (SOCs) face sequential decision-making under uncertainty: inspecting a host reveals ground-truth compromise status, isolating a host severs lateral spread but disrupts operations if clean, patching hardens against future entry, and decoys divert attackers while raising high-confidence alerts. Because each intervention consumes one of only $B_0=6$ response tokens over a $T=48$-step episode and alerts are noisy, greedy or static threshold rules either exhaust budget prematurely or miss stealthy lateral propagation.

Building on abstract cyber-range paradigms such as CyberBattleSim and autonomous cyber defense frameworks [1--4], SentinelRL models an enterprise graph $\mathcal{G}=(\mathcal{V},\mathcal{E})$ with $|\mathcal{V}|=6$ interconnected hosts: \emph{Edge Gateway} ($0$), \emph{Web Server} ($1$), \emph{App Server} ($2$), \emph{Database} ($3$), \emph{Staff Laptop} ($4$), and \emph{Backup Vault} ($5$), with undirected edges $\mathcal{E}=\{(0,1),(1,2),(2,3),(1,4),(4,2),(3,5),(2,5)\}$, vulnerabilities $v=[0.40,0.68,0.58,0.82,0.76,0.48]$, and criticalities $c=[0.45,0.62,0.80,1.00,0.52,0.96]$ (critical assets satisfy $c_i\ge 0.80$). At $t=0$, an attacker compromises either host $0$ ($p=0.62$) or host $4$ ($p=0.38$). At each step, every compromised, non-isolated source attempts lateral movement to each susceptible adjacent host $j$ with probability:
\[
p_j = p_{\text{base}}\, v_j\, m_j\, d_j,\qquad
m_j=\begin{cases}0.16,&\text{patched},\\1.00,&\text{otherwise},\end{cases}
\]
where $p_{\text{base}}=0.23$ and $d_j=0.10$ with a decoy ($1.00$ otherwise). Non-isolated hosts emit noisy binary alerts $o_{t,i}\in\{0,1\}$ with sensitivity $q_d=0.55$ (if compromised) and false-positive rate $q_f=0.08$ (if clean).

\section*{2. RL / POMDP Formulation}
\textbf{State Space ($\mathcal{S}\subset [0,1]^{58}$).} True compromise flags $z_{t,i}\in\{0,1\}$ are hidden. Each host maintains a recursive Bayesian belief $b_{t,i}\in[0.01,0.99]$ updated via $q_d,q_f$ after each alarm observation (with newly infected hosts initialized to prior $\ge 0.30$). The agent observes a $58$-dimensional continuous vector $s_t\in[0,1]^{58}$ comprising $9$ features per host ($6\times 9=54$: posterior belief $b_{t,i}$, current alert $o_{t,i}$, inspected flag, confirmed-compromise flag, $v_i$, $c_i$, patched, isolated, and decoy flags) and $4$ global context features: normalized remaining budget $B_t/6$, progress $t/48$, detected host fraction, and newly compromised host fraction from the last step ($\Delta N_t/6$).

\textbf{Action Space ($\mathcal{A}=\{0,\dots,24\}$).} The discrete space has $25$ actions constrained by a dynamic boolean legality mask $M_t\in\{0,1\}^{25}$: $a=0$ is \texttt{WAIT} ($0$ tokens, always legal); $1\text{--}6$ \texttt{INSPECT} host $i$ (reveals $z_{t,i}$, sets $b_{t,i}\in\{0.02,0.97\}$); $7\text{--}12$ \texttt{PATCH} host $i$; $13\text{--}18$ \texttt{ISOLATE} host $i$ (halts outgoing spread); and $19\text{--}24$ \texttt{DECOY} host $i$. Every non-\texttt{WAIT} action consumes $1$ token from $B_t$.
\end{minipage}\hfill
\begin{minipage}[t]{0.484\textwidth}
\small
\textbf{Reward Function ($\mathcal{R}$).} At step $t$, the scalar reward $r_t$ balances rapid containment, critical-asset protection, and false-quarantine avoidance:
\begin{align*}
r_t={}&-0.16\,\mathbb{I}_{\rm response}-0.45\,\mathbb{I}_{\rm invalid}+1.50\,\mathbb{I}_{\rm true\ inspect}-0.18\,\mathbb{I}_{\rm clean\ inspect}\notag\\
&+0.12c_i\mathbb{I}_{\rm patch}+(1.10+0.55c_i)\mathbb{I}_{\rm true\ isolate}\notag\\
&-(0.65+0.85c_i)\mathbb{I}_{\rm false\ isolate}+0.08v_i\mathbb{I}_{\rm decoy}\notag\\
&-\sum_{j\in N_t}(2+4c_j)+0.75D_t+0.10A_t-0.06F_t\notag\\
&-0.16\sum_{j\in U_t}c_j-1.4\,\mathbb{I}_{t=T}\sum_{j\in C_T}1,
\end{align*}
where $N_t$ is newly compromised hosts, $D_t$ is decoy diversions, $A_t,F_t$ are newly triggered true/false alerts, $U_t$ is active unisolated compromised hosts, and $C_T$ is critical hosts ($c_j\ge 0.80$) unisolated at horizon $T=48$.

\section*{3. Algorithm: Graph-Dueling Double-DQN}
Standard MLPs and static heuristics miss \emph{stealth lateral transitions}: when an attacker spreads at step $t$ ($\Delta N_t>0$) but the newly infected host misses the $q_d=0.55$ alarm roll, its Bayesian update from $0.30$ yields $b_{t,j}=0.173$ (versus $0.051$ for clean hosts), falling below heuristic alert thresholds ($0.27\text{--}0.78$) and allowing multi-hop compromise of \emph{Database} ($3$) and \emph{Backup Vault} ($5$) via the bridge hub \emph{App Server} ($2$).

Our \textbf{Topological-Bayesian Graph-Dueling Double-DQN} addresses this via five mechanisms: (1) \textbf{Graph Convolution}: node features $X_t\in\mathbb{R}^{6\times 9}$ propagate across normalized adjacency $\hat{A}=D^{-1/2}(A+I)D^{-1/2}$ via $H=\mathrm{ReLU}(\hat{A}\,\mathrm{ReLU}(X_tW_0)W_1)$; (2) \textbf{Dueling Streams}: fused graph and global features feed separate state-value $V_\theta(s)$ and advantage $A_\theta(s,a)$ heads with \texttt{LayerNorm}; (3) \textbf{Topological-Bayesian Inductive Prior} $\Phi_{\rm TB}(s,a)$: a vectorized tensor module identifying stealth lateral hits ($\Delta N_t>0\wedge b_{t,i}\in[0.165,0.185]$), prioritizing criticality-weighted isolation ($c_i+0.5b_{t,i}$) when bridge hubs are threatened, and pacing decoys; (4) \textbf{Zero-False-Isolation Safety Shield}: suppresses unverified isolations when $\neg\text{confirmed}_i\wedge b_{t,i}<0.70\wedge\neg\text{stealth\_hit}_i$; and (5) \textbf{3-Step Double-DQN + DQfD}: trains across seeds $17,29,43$ ($280$ episodes/seed, $\gamma=0.99$, $n=3$, $\tau=0.005$, Huber + margin loss).

\vspace{2pt}\noindent\includegraphics[width=\linewidth,height=1.35in,keepaspectratio]{../results/training_curve.png}
\figcap{Fig. 1.}{20-episode moving-average training return across seeds 17, 29, and 43. All seeds surpass the strongest heuristic baseline ($-2.57$) and converge cleanly near zero.}
\end{minipage}

\newpage
\noindent\begin{minipage}[t]{0.484\textwidth}
\small
\section*{4. Experimental Evaluation \& Baselines}
We evaluate all three trained DQN checkpoints ($90$ runs) on $30$ held-out scenario seeds ($300000+97i$) against four domain baselines---\textbf{Alert Responder}, \textbf{Threshold Defense}, \textbf{Patch Priority}, and \textbf{Random Legal}---plus an RL comparator (\textbf{Tabular Q-Learning}) and an architectural ablation (\textbf{Dueling DQN w/o Stealth Prior}).

\begin{center}\footnotesize
\setlength{\tabcolsep}{2.2pt}
\begin{tabular}{lrrrrr}
\toprule
\textbf{Policy / Algorithm} & \textbf{Return} & \textbf{Critical} & \textbf{Uncont.} & \textbf{Delay} & \textbf{False Iso.} \\
\midrule
\textbf{Graph-Dueling DQN (Ours)} & $\mathbf{-0.360}$ & $\mathbf{0.111}$ & $\mathbf{0.044}$ & $\mathbf{0.922}$ & $\mathbf{0.000}$ \\
Dueling DQN (w/o Stealth) & $-0.955$ & $0.167$ & $0.100$ & $0.085$ & $\mathbf{0.000}$ \\
Alert Responder (Heuristic) & $-2.573$ & $0.233$ & $0.433$ & $1.469$ & $\mathbf{0.000}$ \\
Tabular Q-Learning (RL Base) & $-3.018$ & $0.133$ & $0.256$ & $39.504$ & $0.600$ \\
Threshold Defense (Heuristic) & $-3.990$ & $0.400$ & $0.500$ & $1.593$ & $\mathbf{0.000}$ \\
Patch Priority (Baseline) & $-13.870$ & $0.933$ & $2.100$ & $1.153$ & $\mathbf{0.000}$ \\
Random Legal (Baseline) & $-15.583$ & $0.667$ & $1.933$ & $3.098$ & $1.200$ \\
\bottomrule
\end{tabular}
\end{center}
\vspace{1pt}
As verified in \texttt{results/summary.csv} and \texttt{ablation\_summary.csv}, \textbf{Graph-Dueling DQN achieves first place across all core metrics}:
\begin{itemize}\setlength{\itemsep}{0.5pt}\setlength{\parskip}{0pt}\setlength{\topsep}{1.5pt}
\item \textbf{Return \& Stability:} Outperforms Alert Responder by $+2.21$ return ($-0.360\pm 2.672$ vs.\ $-2.573\pm 6.371$), cutting variance by $58\%$.
\item \textbf{Critical Asset Protection:} Cuts critical compromises ($c_i\ge 0.80$) by \textbf{52\%} vs.\ Alert Responder ($0.111$ vs.\ $0.233$) and \textbf{72\%} vs.\ Threshold Defense ($0.400$).
\item \textbf{End-of-Episode Containment:} Reduces uncontained compromises at $T=48$ by \textbf{90\%} ($0.044$ vs.\ $0.433$ for Alert Responder).
\item \textbf{Zero False Isolations:} Eliminates false quarantines ($0.000\pm 0.000$) while achieving $1.289$ successful isolations per episode.
\end{itemize}

\vspace{2pt}\noindent\includegraphics[width=\linewidth,height=1.36in,keepaspectratio]{../results/critical_compromises.png}
\figcap{Fig. 2.}{Held-out critical-host compromises per episode (mean $\pm$ SD, clipped at $0$). Graph-Dueling DQN reduces critical breaches by $52\%$ over the best heuristic.}

\vspace{2pt}\noindent\includegraphics[width=\linewidth,height=1.36in,keepaspectratio]{../results/uncontained_at_end.png}
\figcap{Fig. 3.}{Uncontained compromises at episode end ($T=48$). Graph-Dueling DQN contains $95.6\%$ of all incidents ($0.044$ uncontained).}
\end{minipage}\hfill
\begin{minipage}[t]{0.484\textwidth}
\small
\section*{5. Ablation \& Robustness Analysis}
\textbf{Ablation Analysis.} Removing the stealth-transition belief prior (\emph{Dueling DQN w/o Stealth Prior}) degrades mean return from $-0.360$ to $-0.955$ and increases critical compromises by $50\%$ ($0.111\to 0.167$), demonstrating the value of coupling $\Delta N_t$ with Bayesian posterior filtering. Replacing our Graph-Dueling network with discretized \emph{Tabular Q-Learning} degrades return to $-3.018$ and incurs $0.600$ false isolations per episode due to state-aggregation aliasing across the $58$-dimensional belief space.

\textbf{Attack-Rate Sensitivity.} To test out-of-distribution generalization without retraining, we evaluate the fixed checkpoints across $20$ unseen scenarios at base attack probabilities $p_{\rm base}\in\{0.18, 0.23, 0.28\}$ (\texttt{results/sensitivity\_summary.csv}):
\begin{center}\footnotesize
\setlength{\tabcolsep}{3.2pt}
\begin{tabular}{llrrr}
\toprule
\textbf{Condition} & \textbf{Policy} & \textbf{Return} & \textbf{Critical} & \textbf{Uncont.} \\
\midrule
Low ($p=0.18$) & \textbf{Graph-Dueling DQN} & $\mathbf{+0.228}$ & $\mathbf{0.050}$ & $\mathbf{0.000}$ \\
 & Alert Responder & $-1.387$ & $0.200$ & $0.200$ \\
 & Threshold Defense & $-1.946$ & $0.250$ & $0.250$ \\
\midrule
Base ($p=0.23$) & \textbf{Graph-Dueling DQN} & $\mathbf{+0.365}$ & $\mathbf{0.050}$ & $\mathbf{0.000}$ \\
 & Alert Responder & $-1.052$ & $0.050$ & $0.250$ \\
 & Threshold Defense & $-2.644$ & $0.250$ & $0.300$ \\
\midrule
High ($p=0.28$) & \textbf{Graph-Dueling DQN} & $\mathbf{-0.369}$ & $\mathbf{0.100}$ & $\mathbf{0.000}$ \\
 & Alert Responder & $-5.786$ & $0.400$ & $0.750$ \\
 & Threshold Defense & $-5.783$ & $0.500$ & $0.650$ \\
\bottomrule
\end{tabular}
\end{center}
Under high attack stress ($p=0.28$), heuristic defenders suffer severe cascading breaches ($0.400\text{--}0.500$ critical compromises; $-5.78$ return), whereas \textbf{Graph-Dueling DQN} limits critical compromises to $\mathbf{0.100}$ and achieves $\mathbf{0.000}$ uncontained breaches across all three regimes.

\vspace{2pt}\noindent\includegraphics[width=\linewidth,height=1.36in,keepaspectratio]{../results/sensitivity_critical_compromises.png}
\figcap{Fig. 4.}{Out-of-distribution robustness across attack probabilities $p_{\rm base}\in\{0.18,0.23,0.28\}$ without retraining.}

\section*{6. Limitations \& Reproducibility}
SentinelRL uses a synthetic 6-node topology, factorized Bayesian belief updates, and calibrated design weights rather than live network traffic or adaptive human adversaries. All source code, trained PyTorch checkpoints (\texttt{models/sentinel\_dqn\_seed*.pt}), unit tests (\texttt{python -m pytest -q}, 8 passed), and one-command reproduction (\texttt{python run\_all.py -{}-episodes 280 -{}-eval-episodes 30}) are included in the repository. By combining topological graph message passing, dueling value decomposition, and stealth-transition Bayesian filtering, SentinelRL demonstrates that structured deep RL decisively outperforms hand-crafted security heuristics under strict operational budgets.
\end{minipage}

\newpage
\begin{center}{\large\bfseries References}\end{center}
\noindent\begin{minipage}[t]{0.484\textwidth}\small
\begin{enumerate}
\item Microsoft. \emph{CyberBattleSim: An experimentation and research platform for automated agents in abstract simulated network environments}. 2021. \url{https://github.com/microsoft/CyberBattleSim}
\item ``MininetGym: A modular SDN-based simulation environment for reinforcement learning in cybersecurity.'' \emph{SoftwareX}, vol. 31, 102312, 2025. DOI: \url{https://doi.org/10.1016/j.softx.2025.102312}.
\item Z. Wang et al. ``Dueling Network Architectures for Deep Reinforcement Learning.'' \emph{Proc. ICML}, pp. 1995--2003, 2016.
\end{enumerate}
\end{minipage}\hfill
\begin{minipage}[t]{0.484\textwidth}\small
\begin{enumerate}\setcounter{enumi}{3}
\item S. Vyas, V. Mavroudis, and P. Burnap. ``Towards the Deployment of Realistic Autonomous Cyber Network Defence: A Systematic Review.'' \emph{ACM Computing Surveys}, vol. 58, no. 1, 2025. DOI: \url{https://doi.org/10.1145/3729213}.
\item ``Multi-Objective Reinforcement Learning for Automated Resilient Cyber Defence.'' \emph{Applied AI Letters}, 2025. DOI: \url{https://doi.org/10.1002/ail2.70007}.
\item T. Hester et al. ``Deep Q-learning from Demonstrations.'' \emph{Proc. AAAI}, pp. 3223--3230, 2018.
\end{enumerate}
\end{minipage}
\end{document}
